\section{Supported Scope and Future Work}
The accepted domain is a single hart with trusted kernel/firmware and audited non-privileged single-thread LIFO ELF targets. The tested static and musl dynamic paths follow supported link/return conventions, with the loader continuation represented as abandoned. The compiler fixtures disable sibling-call optimization where required by that convention. Generic language-runtime control transfer is a separate operating model.

Two physical pools provide two concurrent monitored slots. Current software retains 12 pause diagnostics per session; the source-reviewed next-pause path disables monitoring and preserves incomplete output. Hardware stress beyond that diagnostic bound was not newly executed. Full-tree inputs stay within 256 frames; deeper folding has no sticky count after full unwind. Classification has 63 rows and 16 FIFO histories per bank, with warmup/eviction/drop states visible in results.

The classification scope covers retired integer 1/2/4/8-byte loads/stores. Additional floating-point, vector, atomic/LRSC and walker observations require new rules and fixtures. Native HPM profiles are the named DTLB request and Cache request-message metrics, with interval attribution and one fixed profile per firmware image. Automatic DTLB combinations, general multi-thread/shared-mm/vfork cases, custom signal/non-local unwinds and task migration remain future validation domains.

Within this declared scope, the acceptance checklist closes the implementation, lifecycle, export, correctness and descriptive performance requirements. Future work can extend the scope deliberately: continuous diagnostic streaming for long sessions, configurable pool capacity, multi-hart counter epochs, control-flow reconciliation, more events, full-capacity long stress, richer symbolization and non-console transfer. The measured call cost motivates a fast-path investigation before broader deployment. Reliable synthesis and a matched no-extension model would complement the present functional/cycle evidence.
